\documentclass[10pt,a4paper]{article}
\usepackage[a4paper,left=1.35cm,right=1.35cm,top=1.15cm,bottom=1.1cm]{geometry}
\usepackage{fontspec}
\usepackage{xcolor}
\usepackage{microtype}
\usepackage{enumitem}
\usepackage{tabularx}
\usepackage{titlesec}
\usepackage[hidelinks]{hyperref}
\usepackage{parskip}
\usepackage{setspace}

\setmainfont{Segoe UI}
\newfontfamily\brandfont{Trebuchet MS}

\definecolor{fieldgreen}{HTML}{174A3A}
\definecolor{sensorlime}{HTML}{A8C64A}
\definecolor{slate}{HTML}{334155}
\definecolor{mist}{HTML}{EAF1ED}
\definecolor{softgray}{HTML}{64748B}

\pagestyle{empty}
\setlength{\parindent}{0pt}
\setlength{\parskip}{2pt}
\setstretch{0.97}
\emergencystretch=2em
\urlstyle{same}

\titleformat{\section}
  {\brandfont\fontsize{10.8}{12}\selectfont\bfseries\color{fieldgreen}}
  {}{0pt}{}[\vspace{-4pt}{\color{sensorlime}\titlerule[1.2pt]}]
\titlespacing*{\section}{0pt}{5pt}{2.5pt}

\newcommand{\brandheader}{%
  \begin{tabularx}{\textwidth}{@{}Xr@{}}
    {\brandfont\fontsize{25}{27}\selectfont\bfseries\color{fieldgreen}\CandidateName} &
    {\small\color{slate}\CandidateLocation\enspace|\enspace\CandidatePhone}\\[-1pt]
    {\brandfont\fontsize{11.5}{13}\selectfont\color{slate}\CandidateTitle} &
    {\small\href{mailto:\CandidateEmail}{\CandidateEmail}}\\[-1pt]
    {\brandfont\small\bfseries\color{sensorlime}\Availability} &
    {\small\href{\LinkedInURL}{\LinkedInText}\enspace|\enspace\href{\GitHubURL}{\GitHubText}}
  \end{tabularx}
  \vspace{3pt}{\color{fieldgreen}\hrule height 1.6pt}\vspace{2pt}{\color{sensorlime}\hrule height .7pt}\vspace{2pt}
}

\newcommand{\role}[4]{%
  \textbf{#1}\hfill{\brandfont\bfseries\color{fieldgreen}#2}\par
  {\itshape\color{slate}#3}\if\relax\detokenize{#4}\relax\else\hfill{\itshape\color{slate}#4}\fi\par
}

\newenvironment{compactitems}
  {\begin{itemize}[leftmargin=1.15em,itemsep=.3pt,topsep=1pt,parsep=0pt,partopsep=0pt]}
  {\end{itemize}}

\newcommand{\tagline}[1]{%
  \colorbox{mist}{\parbox{\dimexpr\linewidth-2\fboxsep}{\small\bfseries\color{fieldgreen}#1}}
}

\newcommand{\portfolioicon}{\textcolor{sensorlime}{\brandfont\bfseries>}}

\newcommand{\CandidateName}{Harsh Sahu}
\newcommand{\CandidateTitle}{Field Application Scientist \& Product Development Manager}
\newcommand{\CandidateLocation}{Wentorf bei Hamburg, Germany}
\newcommand{\CandidatePhone}{+49 176 34306490}
\newcommand{\CandidateEmail}{harshsahu.8085@gmail.com}
\newcommand{\LinkedInText}{linkedin.com/in/atb-harshsahu}
\newcommand{\LinkedInURL}{https://www.linkedin.com/in/atb-harshsahu}
\newcommand{\GitHubText}{github.com/harshest-human}
\newcommand{\GitHubURL}{https://github.com/harshest-human}
\newcommand{\VideoURL}{https://youtu.be/41Y6egqsQDI}
\newcommand{\PaperURL}{https://www.mdpi.com/2073-4433/14/11/1643}
\newcommand{\Availability}{Available February 2027}


\geometry{left=1.85cm,right=1.85cm,top=1.45cm,bottom=1.4cm}
\begin{document}
\fontsize{10}{12.5}\selectfont
\brandheader

\begin{tabularx}{\textwidth}{@{}Xr@{}}
\textbf{Robert Whittington} & 10 August 2026\\
Precision Livestock Advisor --- Europe, LATAM \& Africa &\\
C-Lock Inc. &
\end{tabularx}

\vspace{6pt}
{\brandfont\fontsize{13}{15}\selectfont\bfseries\color{fieldgreen}Proposal: DACH Field Applications \& Product Development}
\vspace{4pt}

Dear Mr Whittington,

I am writing to propose a dual technical-commercial role through which I could help C-Lock expand both its regional support capacity and its product offering across DACH and Northern Europe.

I am a PhD researcher at Leibniz ATB Potsdam and Freie Universität Berlin, specializing in livestock-emissions measurement, mass-flux methodology, NDIR sensing, barn climate and the integration of environmental and animal data. From February 2027, I am seeking to transition into an industry role combining two connected responsibilities.

\textbf{Field applications and regional development.} Based near Hamburg, I could support GreenFeed adoption across Germany, Denmark, the Netherlands and neighboring markets. German institutes often procure major equipment through publicly funded projects, so adoption begins during consortium formation, experimental design and grant preparation. Through the BMEL-funded EmiMod consortium, I collaborate with ATB, the Thünen Institute, LUFA Nord-West, Kiel University and regional stakeholders. I could help investigators specify GreenFeed in proposals, design robust trials, interpret mass-flux data and receive responsive Tier-1 deployment and calibration support after C-Lock-specific training.

\textbf{Complementary product development.} I could also help develop affordable barn-climate and emissions-monitoring tools that connect individual-animal GreenFeed measurements with spatial gas concentrations, ventilation indicators, temperature, humidity and background conditions. This would strengthen whole-system interpretation and create a scalable entry product for farms and research projects unable to begin with reference-grade instrumentation.

This direction builds on ATB's low-cost wireless monitoring work for CO\textsubscript{2}, NH\textsubscript{3}, CH\textsubscript{4} and barn climate. I was not a co-author of the team's 2023 foundation paper. Since 2023, however, I have taken responsibility for continuing this work under Dr.-Ing. David Janke's supervision, with follow-up manuscripts currently in submission. The work has given me practical experience with sensor selection, calibration, interference, stability, field validation and the trade-offs required to turn affordable components into credible scientific tools.

My practical field environment is shown in ATB's \href{https://youtu.be/41Y6egqsQDI}{\textcolor{fieldgreen}{\textbf{video on measuring emissions in dairy buildings}}}; the underlying product-development direction is documented in the \href{https://www.mdpi.com/2073-4433/14/11/1643}{\textcolor{fieldgreen}{\textbf{2023 low-cost wireless sensor study}}}.

I would welcome a conversation about whether this combined field-application and product-development mandate could address C-Lock's immediate regional needs while opening a longer-term barn-level product opportunity.

Kind regards,

\vspace{8pt}
{\brandfont\bfseries\color{fieldgreen}\CandidateName}
\end{document}
